\section{Prefix histories and a plane that compensates area change}\label{bcs:prefix-section}
Reciprocal rectangle maps preserve planar area. Prefix replacements need not: writing a history symbol can change the total prefix length. We now allow arbitrary positive diagonal scalings in the plane, and compensate their area change in the transverse direction. The computational plane itself remains invariant. A logarithmic clock then gives a complete decaying-force construction from ordinary tape instructions.

\begin{theorem}\label{bcs:prefix-theorem}
For fixed computable $\nu>0$, every deterministic one-tape machine and finite input effectively determine a smooth mean-zero force on the unit torus, with zero initial velocity, such that
\[
 \|\partial_t^jD_x^\alpha f(t)\|_\infty\le C_{j,\alpha}(1+t)^{-1-j}
\]
for every $j,\alpha$. The unique global smooth velocity satisfies the same bounds. Its particle from $(1/4,1/2,1/2)$ enters $\{1/2<x_1<1\}$ exactly when the machine halts. One choice of force has pressure zero; projecting that force gives a divergence-free choice with the same velocity and a changed normalized pressure. For a fixed machine, the repeated motion before the clock change is independent of the input; only its initial loading motion changes.
\end{theorem}

\subsection{A record between ordinary tape symbols}
Let $Q$ be the finite state set, $q_{\mathrm{in}}$ its initial state, $Q_h\subseteq Q$ its terminal states, and $\Gamma$ the work alphabet with blank $\square$. Missing instructions are replaced by stationary transitions to an added terminal state, preserving the read letter. A transition into a terminal state is completed; an initially terminal configuration also counts as halting. Add a read-only origin bit to the work alphabet, one only at cell zero and preserved by every write. For empty input, the initial word is a marked blank. This allows recovery of the absolute head position after decoding.

The record symbols retain erased information in the sense of reversible history computation~\cite{Bennett1973}. For each nonterminal $q$, previous displacement $d\in\{-1,0,1\}$, and read letter $a$, create a record $g(q,d,a)$, and let $G$ be the set of these records. Two infinite stacks $L,R$ use the alphabet $\mathcal A=\Gamma\sqcup G$. Labels are $C(q,d)$ and $S_+(q),S_-(q)$. At a checkpoint $C(q,d)$, deleting the records from $L$ gives the tape left of the head in reverse order; deleting them from $R$ gives the tape from the head rightward. The top of $R$ is a work letter. The scan labels transfer leading records between the stacks until the next work letter appears.

A prefix rule $(\ell;\alpha,\beta)\to(\ell';\alpha',\beta')$ replaces the indicated prefixes of $L,R$ and retains their infinite tails. For $\delta(q,a)=(q',b,m)$ and $g=g(q,d,a)$, use
\begin{equation}\label{bcs:prefix-work-rules}
\begin{array}{lll}
 (C(q,d);\epsilon,a)\to(S_+(q');bg,\epsilon),& &m=1,\\
 (C(q,d);\epsilon,a)\to(S_-(q');\epsilon,bg),&&m=-1,\\
 (C(q,d);\epsilon,a)\to(C(q',0);g,b),&&m=0.
\end{array}
\end{equation}
Here $\epsilon$ is the empty word. For every $q$, every $g\in G$ and every $c\in\Gamma$, also use
\begin{equation}\label{bcs:prefix-scan-rules}
\begin{array}{ll}
 (S_+(q);\epsilon,g)\to(S_+(q);g,\epsilon),&
 (S_+(q);\epsilon,c)\to(C(q,1);\epsilon,c),\\
 (S_-(q);g,\epsilon)\to(S_-(q);\epsilon,g),&
 (S_-(q);c,\epsilon)\to(C(q,-1);\epsilon,c).
\end{array}
\end{equation}
The terminal labels are exactly the checkpoints $C(q,d)$ with $q\in Q_h$;
there are no rules out of them. The scan labels $S_\pm(q)$ remain
nonterminal even for $q\in Q_h$, so a pending move finishes before its
terminal checkpoint is reached. Initialize at $C(q_{\mathrm{in}},0)$ with
left stack $\square^\infty$ and right stack $\omega\square^\infty$, where
$\omega$ includes the origin bit.

\begin{lemma}\label{bcs:prefix-symbolic}
The full domain cylinders of these rules are pairwise disjoint, and their full image cylinders are pairwise disjoint. Initialized checkpoints recover the ordinary machine configurations. Each ordinary transition takes a positive finite number of prefix steps, and a terminal checkpoint occurs exactly on halting.
\end{lemma}
\begin{proof}
Rules with a common source label test incompatible top letters. For images entering $S_+(q)$, an ordinary transition begins the left stack with a work letter followed by its distinct source record, whereas a scan loop begins it with a record. These distinguish all incoming cylinders. The same argument uses the right stack for $S_-(q)$. Into $C(q,1)$ or $C(q,-1)$, the final scan writes a distinct first work letter on $R$; into $C(q,0)$, the source record at the start of $L$ distinguishes the rule. Labels separate the remaining cases. Thus these assertions hold for arbitrary tails, without an initialization assumption.

For a right move, \eqref{bcs:prefix-work-rules} removes the read letter from $R$ and puts the written letter and its record on $L$. The positive scan transfers leading records from $R$ until the next work letter is exposed. Deleting records gives exactly the rightward tape update. A left move first replaces the read letter by $bg$ on $R$; its scan transfers records from $L$ and then its first work letter to $R$, giving the leftward update. A stationary move only changes the scanned work letter after deletion. Each transition adds one record, so there are finitely many records at every reached checkpoint and every scan is finite. The origin bit survives deletion and gives absolute alignment. A scan associated to a transition into $Q_h$ finishes before the terminal checkpoint is declared. This proves the assertions, including initial halting.
\end{proof}

\subsection{A planar motion for arbitrary diagonal rectangle maps}
As in Moore's positional representation of symbolic shifts~\cite{Moore1991}, assign the ordered letters of $\mathcal A$ the digits $1,3,\ldots,2M-1$ in base $B=2M+1$, where $M=|\mathcal A|$. Set
\[
 \gamma(w)=\sum_{j\ge1}D(w_j)B^{-j},\quad
 c(\alpha)=\sum_{j=1}^{|\alpha|}D(\alpha_j)B^{-j},\quad
 I(\alpha)=[c(\alpha),c(\alpha)+B^{-|\alpha|}].
\]
The identity $\gamma(\alpha w)=c(\alpha)+B^{-|\alpha|}\gamma(w)$ and the gaps between digits prove unique decoding. Child intervals lie strictly inside the parent, so incompatible prefix cylinders have separated intervals even as closed sets.

Put $\Omega=(0,1)^2$ and $\Omega_-=(0,1/2)\times(0,1)$. If the labels number $N$, give label $\ell_i$ a square $K_i$ of half-side $1/[16(N+1)]$, centered at height $i/(N+1)$ and first coordinate $3/4$ for terminal labels, $1/4$ otherwise. Let $\Phi_i$ be the positive axis-preserving affine map of $[0,1]^2$ onto that square. The encoded state is
$\Phi_i(\gamma(L),\gamma(R))$. A rule with prefixes $\alpha,\beta$ and $\alpha',\beta'$ has rectangles
\[
 D_i=\Phi_\ell(I(\alpha)\times I(\beta)),\qquad
 E_i=\Phi_{\ell'}(I(\alpha')\times I(\beta')).
\]
Both families are separately disjoint by Lemma~\ref{bcs:prefix-symbolic}. Each source lies in $\Omega_-$, and each target lies entirely in one open half of $\Omega$. The positive diagonal affine map $D_i\to E_i$ implements the prefix rule exactly. Its determinant is allowed to differ from one.

The shrink--move--expand interpolation also appears in the proof of~\cite[Proposition~5.1]{CardonaMirandaPeraltaSalasPresas2021}; here we specify rational routes that preserve the required observation half-plane.

\begin{lemma}\label{bcs:prefix-planar}
Let $D_i,E_i$, $1\le i\le m$, be rational axis-aligned closed rectangles of positive side lengths. Suppose the $D_i$ are pairwise disjoint in $\Omega_-=(0,1/2)\times(0,1)$, the $E_i$ are pairwise disjoint in $\Omega=(0,1)^2$, and each $E_i$ lies entirely in one open half of $\Omega$. There is an effective smooth field $v_*(s,Y)$ supported in $(0,1)\times\Omega$ whose time-one map on each full $D_i$ is its positive diagonal affine identification with $E_i$. If $E_i\subset\Omega_-$, the whole trajectory from $D_i$ remains in $\Omega_-$. No incompressibility is asserted in this planar lemma.
\end{lemma}
\begin{proof}
For an empty list take zero. Otherwise choose positive rational enlargement margins separately for the source and target families, preserving their disjointness and their open-half containment. Let $p_i,q_i$ be their centers, and $r_i^D,r_i^E$ their half-side vectors. Write $y_0=\min_i\{p_{i,2},q_{i,2}\}$ and choose temporary centers
\[
 k_i=(1/8,\ i y_0/[2(m+1)]),\qquad 1\le i\le m.
\]
These are mutually distinct and lie below all source and target centers. Move $p_i$ to $k_i$ in increasing order, avoiding the centers $k_j$ with $j<i$ and $p_j$ with $j>i$. Then move $k_i$ to $q_i$ in increasing order, avoiding $q_j$ with $j<i$ and $k_j$ with $j>i$. Each first route stays in $\Omega_-$; a second route stays there if its target does, and otherwise stays in $\Omega$.

Here is a finite rational rule for such a route. In its allowed rectangle $(0,b)\times(0,1)$, choose a point $w$ at height $\min(p_2,q_2)/2$, off every line from either endpoint to a forbidden center. No horizontal forbidden line has this height. Each nonhorizontal line has one computable rational intersection with it. Choose $w_1$ positive and smaller than $b$ and every positive intersection coordinate, for example half their positive minimum with $b$. The two segments $p$--$w$--$q$ stay in the allowed rectangle and avoid the forbidden centers.

Choose a common rational $\varepsilon>0$ equal to $1/20$ times the minimum of the rectangle half-sides, segment-endpoint boundary distances, and, for each segment $a$--$b$ and forbidden center $o$, $|\det(b-a,o-a)|/2$; omit empty collections. Every included number is positive. Since segment lengths are less than two, the determinant quantity is a lower bound for the Euclidean distance to its line. The moving square of half-side $2\varepsilon$ therefore misses every stationary square of half-side $\varepsilon$: an intersection would put their centers within $3\sqrt2\varepsilon$, contradicting this clearance. It also remains inside its allowed rectangle.

Shrink each $D_i$ about $p_i$ to a square of half-side $\varepsilon$. To implement this, let $\zeta$ be a smooth step constant near the endpoints of $[0,1]$, put
\[
 A_i(\tau)=(1-\zeta(\tau))I+
       \zeta(\tau)\operatorname{diag}((\varepsilon,\varepsilon)/r_i^D),
\]
and multiply $A_i'A_i^{-1}(Y-p_i)$ by a cutoff equal to one near $D_i$ and supported in its chosen enlargement. The cutoffs are disjoint, all diagonal entries are positive, and the exact path is $p_i+A_i(\tau)(Y-p_i)$. The matrices shrink coordinatewise, so these paths stay in $D_i$.

For each scheduled segment $a$--$b$, put $c(\tau)=a+\zeta(\tau)(b-a)$ and use the translation field $c'(\tau)$ times a cutoff which is one on a neighborhood of the moving square and supported in the concentric square of half-side $2\varepsilon$. This translates that square exactly and vanishes on all held squares. After every small square has reached its target center, reverse the shrinking operation at each $q_i$ with matrices $(1-\zeta)I+\zeta\operatorname{diag}(r_i^E/(\varepsilon,\varepsilon))$. Its paths remain in $E_i$.

Each field vanishes near its phase endpoints. Rescale the finite list of $J$ phases $V_0,\ldots,V_{J-1}$ by $v_*(s)=J\sum_jV_j(Js-j)$ to fit them in one time unit. The resulting map sends $p_i+\xi$ to $q_i+\operatorname{diag}(r_i^E/r_i^D)\xi$, as required. Every stage obeys the indicated half-plane restriction. All cutoffs, paths, inverse diagonal lower bounds and derivative estimates are finite effective data.
\end{proof}

\subsection{Compensating the planar divergence}
The rational initial code is denoted by $Z_{\mathrm{in}}$. Its rationality follows from the finite input and the blank tail $D(\square)/(B-1)$. Let $p_*=(1/4,1/2)$ and load it along
$c(s)=p_*+\zeta(s)(Z_{\mathrm{in}}-p_*)$.
A translating bump supported in a sufficiently thin tube around this segment gives a field $v_{\mathrm{in}}$ with precisely that trajectory. For instance, take the half-side of the moving square to be one quarter of the smaller endpoint distance from $\partial\Omega$, and take the same enlargement margin. Convexity keeps the support in $\Omega$. A nonterminal loader stays in $\Omega_-$; a terminal loader reaches the right half.

Extend $v_*,v_{\mathrm{in}}$ by zero in time and periodically in space, and define
\[
 v(s,Y)=v_{\mathrm{in}}(s,Y)+\sum_{n\ge1}v_*(s-n,Y).
\]
All mixed derivatives are bounded because only one repeated phase template and one loader occur. At logical time $1+n$, the trajectory from $p_*$ is the code after $n$ prefix steps, up to a terminal visit. The rectangle lemma proves absence of a right-half visit during every nonterminal transition as well as at its endpoints. Together with the finite-scan invariant, this gives the all-time halting equivalence.

Choose a periodic smooth function $\eta(z)$ supported near $1/2$ which equals $z-1/2$ on a neighborhood of $1/2$. Define
\begin{equation}\label{bcs:prefix-lift}
 W(s,Y,z)=\bigl(\eta'(z)v(s,Y),-\eta(z)\operatorname{div}_Yv(s,Y)\bigr).
\end{equation}
The horizontal divergence and vertical derivative cancel. The spatial mean is zero: $\int\eta'=0$, and the integral of a periodic planar divergence is zero. At $z=1/2$, the vertical component vanishes and the horizontal component is exactly $v$. Thus this plane carries the entire prescribed computation despite the planar area changes. The lift has bounded mixed derivatives and vanishes at logical time zero.

Apply Lemma~\ref{bcs:clock} with $g(t)=\log(1+t)$. The particle from $(p_*,1/2)$ is $(Y(g(t)),1/2)$, so its prefix-step observations occur at $t=e^{1+n}-1$. Every finite logical time occurs at finite physical time, and the event equivalence is preserved. The lemma gives the exact derivative decay for the velocity and its residual, which has zero mean and pressure zero. Lemma~\ref{lem:b-comparison} gives uniqueness on the torus.

For clarity, the optional solenoidal force changes pressure. Given the residual $f$, let $\phi$ be the zero-mean solution of $\Delta\phi=\operatorname{div}f$ and replace $(f,p)$ by $(f-\nabla\phi,-\phi)$. This leaves $-\nabla p+f$ and therefore the velocity unchanged. In Fourier variables,
\[
 \phi(t,x)=-\sum_{k\ne0}\frac{i\,k\cdot\widehat f_k(t)}{2\pi|k|^2}e^{2\pi i k\cdot x}.
\]
Integration by parts in a coordinate of maximal $|k_l|$ gives, for every $m$,
$|\partial_t^j\widehat f_k(t)|\le C_{j,m}(1+t)^{-1-j}|k|_\infty^{-m}$.
There are $24r^2+2$ lattice points with $|k|_\infty=r$. Choosing $m$ large proves uniform convergence after any fixed derivatives, retains the decay, and gives effective tails. Finite Fourier coefficients are effectively integrated using derivative bounds. Termwise differentiation gives the Poisson equation, with zero mode zero; uniqueness of Fourier coefficients on continuous periodic functions identifies it pointwise. Thus the projected force is effective, mean zero and divergence free. This is the classical Helmholtz--Leray projection, with its pressure correction explicit.

Finally all infinite sums used to prescribe the force are evaluated from finite supersets of possibly active translates, obtained from a coarse enclosure of the time and position. The elementary flat-step derivatives have effective bounds near their seams. Neither construction nor evaluation uses a configuration reached after initialization. This completes Theorem~\ref{bcs:prefix-theorem}.

\paragraph{Other onto clocks.}
The field $W$ in \eqref{bcs:prefix-lift}, including its loader, also satisfies Proposition~\ref{bcs:power-clock}. Using $(1+t)^{1/3}-1$ or $(1+t)^{1/4}-1$ therefore gives separate force choices with the same fixed label and detector. Every mixed derivative is respectively $O((1+t)^{-2/3})$ or $O((1+t)^{-3/4})$. After the loader, the code after $n$ prefix steps is sampled at $(2+n)^3-1$ or $(2+n)^4-1$. These assertions use the same full-domain instruction motion and the same all-time path bound.
